% GENERATED FILE. Edit canonical lessons, then run npm run generate:books.
% canonical-source-hash: fnv1a64:ae2428e41bb822ab
% canonical-lessons: RU-C30-hair, RU-C30-finger, RU-C30-neck, RU-C30-shoulder, RU-C30-knee

\chapter{Hair, Finger, Neck, Shoulder, Knee}
\label{ch:ru-hair-finger-neck-shoulder-knee}
\hlchaptermodality{30}

\begin{chapteropening}
\textbf{By the end of this chapter:} \emph{I can name hair, a finger, the neck, a shoulder and a knee.}

Complete the last lesson of chapter 30: i can name hair, a finger, the neck, a shoulder and a knee.
\end{chapteropening}

\section[vólosy]{\ru{волосы} (vólosy) — hair}

\label{lesson:RU-C30-hair}



\begin{quote}
Before the new one: say the Russian for the belly, then the Russian for a tooth.
\end{quote}

\subsection*{You'll want to know: \ru{волосы}}
\textbf{\ru{волосы}} (\emph{vólosy}) — \textquotedblleft{}hair\textquotedblright{}. Plural, like English \textquotedblleft{}locks\textquotedblright{}.

\begin{grammarlens}[title={What you've built}]
One more word for hair, finger, neck, shoulder, knee.
\end{grammarlens}

\subsection*{Your turn}
\begin{itemize}
\raggedright
  \item \emph{Say it:} \emph{vólosy}
  \item \emph{Say it:} \emph{vólosy}, once more
  \item \emph{Say it:} say \emph{vólosy}
  \item \emph{Recall:} say the Russian for the belly, then the Russian for a tooth, then say \emph{vólosy} again
\end{itemize}

\begin{tcolorbox}[breakable,colback=teal!4,colframe=teal!35!black,title={Before you move on}]
What does \ru{волосы} mean? (\textquotedblleft{}Hair\textquotedblright{}.) Say it once more.
\end{tcolorbox}
\section[pálets]{\ru{палец} (pálets) — a finger}

\label{lesson:RU-C30-finger}



\begin{quote}
Before the new one: say the Russian for a tooth, then the Russian for hair.
\end{quote}

\subsection*{You'll want to know: \ru{палец}}
\textbf{\ru{палец}} (\emph{pálets}) — \textquotedblleft{}a finger\textquotedblright{}.

\begin{grammarlens}[title={What you've built}]
One more word for hair, finger, neck, shoulder, knee.
\end{grammarlens}

\subsection*{Your turn}
\begin{itemize}
\raggedright
  \item \emph{Say it:} \emph{pálets}
  \item \emph{Say it:} \emph{pálets}, once more
  \item \emph{Say it:} say \emph{pálets}
  \item \emph{Recall:} say the Russian for a tooth, then the Russian for hair, then say \emph{pálets} again
\end{itemize}

\begin{tcolorbox}[breakable,colback=teal!4,colframe=teal!35!black,title={Before you move on}]
What does \ru{палец} mean? (\textquotedblleft{}A finger\textquotedblright{}.) Say it once more.
\end{tcolorbox}
\section[shéya]{\ru{шея} (shéya) — the neck}

\label{lesson:RU-C30-neck}



\begin{quote}
Before the new one: say the Russian for hair, then the Russian for a finger.
\end{quote}

\subsection*{You'll want to know: \ru{шея}}
\textbf{\ru{шея}} (\emph{shéya}) — \textquotedblleft{}the neck\textquotedblright{}.

\begin{grammarlens}[title={What you've built}]
One more word for hair, finger, neck, shoulder, knee.
\end{grammarlens}

\subsection*{Your turn}
\begin{itemize}
\raggedright
  \item \emph{Say it:} \emph{shéya}
  \item \emph{Say it:} \emph{shéya}, once more
  \item \emph{Say it:} say \emph{shéya}
  \item \emph{Recall:} say the Russian for hair, then the Russian for a finger, then say \emph{shéya} again
\end{itemize}

\begin{tcolorbox}[breakable,colback=teal!4,colframe=teal!35!black,title={Before you move on}]
What does \ru{шея} mean? (\textquotedblleft{}The neck\textquotedblright{}.) Say it once more.
\end{tcolorbox}
\section[plechó]{\ru{плечо} (plechó) — a shoulder}

\label{lesson:RU-C30-shoulder}



\begin{quote}
Before the new one: say the Russian for a finger, then the Russian for the neck.
\end{quote}

\subsection*{You'll want to know: \ru{плечо}}
\textbf{\ru{плечо}} (\emph{plechó}) — \textquotedblleft{}a shoulder\textquotedblright{}.

\begin{grammarlens}[title={What you've built}]
One more word for hair, finger, neck, shoulder, knee.
\end{grammarlens}

\subsection*{Your turn}
\begin{itemize}
\raggedright
  \item \emph{Say it:} \emph{plechó}
  \item \emph{Say it:} \emph{plechó}, once more
  \item \emph{Say it:} say \emph{plechó}
  \item \emph{Recall:} say the Russian for a finger, then the Russian for the neck, then say \emph{plechó} again
\end{itemize}

\begin{tcolorbox}[breakable,colback=teal!4,colframe=teal!35!black,title={Before you move on}]
What does \ru{плечо} mean? (\textquotedblleft{}A shoulder\textquotedblright{}.) Say it once more.
\end{tcolorbox}
\section[koléno]{\ru{колено} (koléno) — a knee}

\label{lesson:RU-C30-knee}



\begin{quote}
Before the new one: say the Russian for the neck, then the Russian for a shoulder.
\end{quote}

\subsection*{You'll want to know: \ru{колено}}
\textbf{\ru{колено}} (\emph{koléno}) — \textquotedblleft{}a knee\textquotedblright{}.

\begin{grammarlens}[title={What you've built}]
One more word for hair, finger, neck, shoulder, knee.
\end{grammarlens}

\subsection*{Your turn}
\begin{itemize}
\raggedright
  \item \emph{Say it:} \emph{koléno}
  \item \emph{Say it:} \emph{koléno}, once more
  \item \emph{Say it:} say \emph{koléno}
  \item \emph{Recall:} say the Russian for the neck, then the Russian for a shoulder, then say \emph{koléno} again
\end{itemize}

\begin{tcolorbox}[breakable,colback=teal!4,colframe=teal!35!black,title={Before you move on}]
What does \ru{колено} mean? (\textquotedblleft{}A knee\textquotedblright{}.) Say it once more.
\end{tcolorbox}
